\chapter{Tubería de presión, transitorios y estructuras}
\section{Material, diámetro y presión estática}
Se propone tubería de acero soldado, diámetro interior nominal de 1000 mm y espesor nominal de 16 mm, con recubrimiento exterior e interior y protección según agresividad del medio. El diámetro exterior nominal es 1032 mm. Se especifica un límite de fluencia mínimo de 250 MPa como requisito de compra, sin atribuirlo a una tubería existente. Se consideran tolerancia negativa de espesor de 12.5\% y pérdida futura por corrosión de 1 mm: espesor resistente al final de vida de \dato{EspEf} mm. Estos márgenes son hipótesis que deberán sustituirse por especificación de fabricación y evaluación de corrosión.

La presión estática asociada al salto de 45.3 m es aproximadamente 4.44 bar manométricos en el extremo bajo. En régimen, las pérdidas reducen la presión disponible, pero el cierre puede generar una sobrepresión mucho mayor. Una clase nominal no se selecciona únicamente a partir del salto estático. Alternativas en material plástico requerirían verificar temperatura, fluencia, uniones, anclajes, instalación, golpe de ariete y disponibilidad comercial a gran diámetro; no se les atribuye resistencia equivalente al acero sin cálculo.

\section{Golpe de ariete y límites del modelo}
Se desarrolla un modelo unidimensional por el método de las características (MOC): embalse de carga de nivel fijo, tubería uniforme, fricción Darcy, y válvula equivalente con \(Q=k\tau\sqrt{H}\). Se usa paso espacial uniforme con 100 segmentos y paso temporal que cumple Courant 1. La celeridad se calcula con compresibilidad del agua y elasticidad del acero. Se analizan cierres de 0, 1, 4 y 8 s tanto para pared nueva mínima de 14 mm como para pared final de 13 mm. El diámetro interior aumenta coherentemente al reducir pared; se incrementa también rugosidad a 0.15 mm al final de vida.

El caudal inicial es 1.90 \(\mathrm{m^3/s}\), mayor que el caudal nominal de operación, para estudiar el máximo hidráulico. Se incluye un perfil ideal de descenso continuo; los puntos altos reales y el cierre efectivo de una Banki pueden modificar las ondas. La relación de Joukowsky \(\Delta H=a\Delta v/g\) se utiliza como comprobación del primer salto para un cierre instantáneo sin fricción, no como sustituto del análisis temporal.
\input{generado/moc}
Se detiene la integración al detectar presión manométrica inferior a una referencia de --6.8 m de columna de agua, representativa del límite atmosférico--vapor a la altitud del sitio. El valor es un supuesto conservador, no una medición barométrica. Después de ese instante el modelo monofásico pierde validez: no se publican presiones algebraicas físicamente imposibles ni se certifica el pico posterior al colapso de cavidades. El cierre rápido presenta este problema.

Para cierre controlado de ocho segundos se obtiene una máxima de estudio cercana a \dato{PNslow} bar. La primera máxima válida de cierre brusco alcanza aproximadamente \dato{PNfast} bar. Se estudian válvulas y accesorios \textbf{PN40 como candidatos}, pero esa clase no demuestra por sí sola seguridad ante separación de columna, vacío, fatiga o colapso de cavidades. Se exige validar con el perfil real, curva de cierre, condiciones de admisión, regulación y modelo bifásico cuando corresponda. El ELC debe tomar el rechazo rápido mientras el gobierno cierra lentamente. Debe evitarse que una orden eléctrica provoque cierre hidráulico instantáneo.
\begin{figure}[htbp]\centering\includegraphics[width=.95\linewidth]{figuras/transitorio-r03.png}\caption{Resultados válidos del MOC y envolvente del cierre controlado. Las trazas truncadas indican pérdida de validez por presión baja.}\end{figure}

\section{Espesor resistente y comprobación de membrana}
Se evalúa esfuerzo circunferencial \(\sigma_\theta=pD/(2t_e)\), longitudinal de membrana \(\sigma_z\simeq\sigma_\theta/2\) y equivalente de von Mises. Con el estado de pared y la presión compatibles, resultan \(\sigma_\theta=\dato{SigmaHoop}\) MPa y \(\sigma_{VM}=\dato{SigmaVM}\) MPa. El admisible de estudio es \(250\times0.85/1.5=\dato{SigmaAllow}\) MPa, usando eficiencia de junta 0.85 y factor 1.5 \textbf{supuestos}, no valores atribuidos a una cláusula no consultada.

La comprobación de membrana es favorable bajo esos supuestos, pero no cubre flexión local de silletas, ovalización, pandeo por vacío, concentraciones en derivaciones, fatiga ni estabilidad sísmica. ASCE MOP 79, segunda edición 2012 (N26), identifica criterios complementarios para tuberías de acero; se consultó su ficha y estructura, no el texto completo. USBR EM03 (1966, N20) se utiliza como referencia histórica de mecanismos, sin adoptar códigos antiguos como obligación peruana.

\section{Apoyos, anclajes y acciones térmicas}
Se comparan vanos de 4, 6 y 8 m mediante una viga ideal bajo peso de acero y agua. Para metrados se usa separación preliminar de 4 m, que produce 89 apoyos incluyendo extremos. La masa nominal de acero es aproximadamente \dato{MassPipe} t antes de accesorios. El cálculo global no dimensiona una silleta: deben verificarse contacto, ángulo de apoyo, refuerzos anulares, reacciones, fundación, deslizamiento y posibilidad de movimiento térmico.
\input{generado/apoyos}

Para \(\Delta T=40\ ^\circ\mathrm{C}\), la expansión libre de los 350 m es aproximadamente \dato{Expansion} mm. Se requiere una disposición coordinada de puntos fijos, apoyos deslizantes y juntas, sin permitir fuerzas térmicas incompatibles con los anclajes. En válvulas, codos y cambios de pendiente se evalúan empujes y fuerzas dinámicas; los anclajes no pueden ser simples dados sin verificación.

Una equivalencia de peso--fricción para un codo de estudio de 45 grados, usando \(\mu=0.6\), peso específico de concreto de 24 \(\mathrm{kN/m^3}\) y factor 1.5, arroja un volumen de gravedad cercano a \dato{AnchorVol} \(\mathrm{m^3}\) ante el primer pico brusco válido. Este resultado muestra la magnitud del problema; \textbf{no es un bloque de concreto diseñado}. Se considera estudiar anclaje a roca competente, con caracterización geotécnica, extracción, arrancamiento, estabilidad global y pruebas. Se presupuestan conjuntos provisionales y no un volumen ejecutivo inventado.

\section{Estructuras, geotecnia y sismo}
El diseño académico adopta la E.030 modificada por RM 183-2026-VIVIENDA y considera la disposición transitoria de RM 217-2026. El anexo distrital sitúa Limbani en zona 2, \(Z=0.25\). La categoría A2 comprende instalaciones de generación y transformación, con \(U=1.5\). No se pueden completar espectro, fuerzas y derivas sin perfil de suelo, sistema estructural, períodos y estudio de mecánica de suelos. E.050 (2018) guía la necesidad y contenido del estudio geotécnico; E.060 y demás normas estructurales aplicables se verifican antes de diseñar armado y fundaciones. El acero de refuerzo usado en presupuesto es un indicador de metrado, no un cálculo estructural.

Se proponen válvula principal de aislamiento, órganos de admisión por grupo, purgas, drenes, medición de presión y dispositivos de aire donde el perfil lo requiera. Su tamaño, tipo, capacidad y tiempos de actuación deben coordinarse con el MOC. No se adopta una válvula de aire genérica como solución automática al golpe de ariete. El plan de inspección incluirá soldaduras, espesores, recubrimiento, estanqueidad, instrumentación y pruebas de presión definidas por la especificación aplicable.
